% Both panels and their common caption form one full-width figure.
% Shared unchanged by the English and Chinese prose editions.
\par\medskip
\noindent\begin{minipage}{\linewidth}
\captionsetup{type=figure}
\centering
\begin{minipage}[t]{0.48\linewidth}
\vspace{0pt}
\centering
{\small\sffamily\bfseries (a) State relations}\par\medskip
% Shared vector source for the English and Chinese prose editions.
\begin{tikzpicture}[
  summary/.style={draw=rulegray,rounded corners=2pt,fill=boxgray!45,
    text width=2.45cm,minimum height=1.05cm,inner sep=4pt,
    align=center,font=\small},
  arrow/.style={-{Stealth[length=2mm]},semithick},
  label/.style={font=\small,fill=white,inner sep=2pt,align=center}]
\node[summary,draw=origblue] (original) at (0,2.2)
  {$A_0(S),\ A_0(T)$\\[3pt]$(t_p,ss_p)$};
\node[summary,draw=auxpurple] (maintained) at (0,0)
  {$A(S),\ A(T)$\\[3pt]$(t_p,ss_p,\aux{n_p})$};
\node[summary,draw=fixgreen] (merged) at (3.65,0)
  {$A(U)$\\[3pt]$(t_U,ss_U,\aux{n_U})$};
\node[summary] (reference) at (0,-2.6)
  {$B(S),\ B(T)$\\[3pt]$(t_p,q_p,n_p)$};
\node[summary] (combined) at (3.65,-2.6)
  {$B(U)$\\[3pt]$(t_U,q_U,n_U)$};
\draw[arrow,draw=auxpurple] (original) --
  node[label,right,text=auxpurple] {add $n_p=\sum_{j\in p}1$\\\S\ref{sec:selection}--\ref{sec:synthesis}} (maintained);
\draw[arrow,draw=fixgreen] (maintained) --
  node[label,above,text=fixgreen] {$\Mrg$}
  node[label,below=6mm,text=fixgreen] {\S\ref{sec:correctness}--\ref{sec:rebase}} (merged);
\draw[arrow] (maintained) -- node[label,left] {$H$\\\S\ref{sec:recovery}} (reference);
\draw[arrow] (reference) -- node[label,above] {$\oplus$}
  node[label,below] {\S\ref{sec:correctness}} (combined);
\draw[arrow] (combined) -- node[label,right] {$G$\\\S\ref{sec:atoms}} (merged);
\end{tikzpicture}

\end{minipage}\hfill
\begin{minipage}[t]{0.48\linewidth}
\vspace{0pt}
\small
\centerline{\sffamily\bfseries (b) Derivation}
\medskip
On each part, let $t=\sum_jx_j$, $q=\sum_jx_j^2$ and $n=\sum_j1$.
Expanding the source definition gives
\[
G(t,q,n)=(t,q-t^2/n,n).
\]
From $(t,ss)$, recovery leaves $q=ss+t^2/n$ dependent on the unknown $n$.
Adding $n$ gives the complete recovery
\[
H(t,ss,n)=(t,ss+t^2/n,n).
\]
Since $n>0$, both maps are defined and $H\circ G=\mathrm{id}$.
The reference join $\oplus$ adds coordinates.
The upper arrow follows by composition:
\[
\Mrg(a,b)=G\bigl(H(a)\oplus H(b)\bigr).
\]
With $t_U=t_S+t_T$ and $n_U=n_S+n_T$, rewriting yields the centred form
\[
\begin{aligned}
ss_U={}&ss_S+ss_T\\
&+\frac{n_Sn_T}{n_U}
  \left(\frac{t_S}{n_S}-\frac{t_T}{n_T}\right)^2.
\end{aligned}
\]
The emitted join evaluates this correction without materialising $q$.
\end{minipage}
\caption{Variance: auxiliary states, recovery and join derivation.
Here $p\in\{S,T\}$ and $U=S\dunion T$ for disjoint non-empty parts.
Arrow annotations identify the subsections that develop each step:
feature construction and synthesis (\S\ref{sec:selection}--\ref{sec:synthesis}), information equivalence (\S\ref{sec:recovery}),
join derivation (\S\ref{sec:correctness}), and numerical formulation (\S\ref{sec:rebase}).}
\label{fig:recovery-square}
\end{minipage}\par\medskip
